% titel: Blatt 7, A2: Charakterisierung der Exponentialfunktion
% skript: Kap. 4 (Differentialrechnung)

% Teil I: Werkzeuge, kein Lösungsweg.

Eine differenzierbare Funktion, deren Ableitung ein Vielfaches ihrer selbst
ist, hat eine ganz bestimmte Gestalt. Der übliche Beweisweg nutzt eine
Hilfsfunktion und die Tatsache, dass eine Funktion mit Ableitung null konstant
ist.

\begin{werkzeugkasten}[Ableitungsregeln mit der Exponentialfunktion]
Für festes $a\in\mathbb R$ ist $\mathrm e^{ax}$ nach $x$ abgeleitet
$a\mathrm e^{ax}$:
\begin{equation}
  \dv{}{x}\mathrm e^{ax}=a\,\mathrm e^{ax}.
\end{equation}
Zusammen mit der Produktregel $(uv)'=u'v+uv'$ lässt sich daraus die Ableitung
von Produkten der Form $f(x)\mathrm e^{ax}$ oder $f(x)\mathrm e^{-ax}$
berechnen. Entscheidend ist das Vorzeichen im Exponenten, denn die innere
Ableitung von $-ax$ ist $-a$.

\textit{Skript:} Kap. 4, Ableitungsregeln und Exponentialfunktion.
\end{werkzeugkasten}

\begin{werkzeugkasten}[Konstante Funktion und Hauptsatz]
Ist eine differenzierbare Funktion $g$ auf einem Intervall in allen Punkten
mit $g'=0$ versehen, so ist $g$ dort konstant. Das ist die Brücke zwischen
der Differentialrechnung und der Aussage, dass die Gestalt von $f$ durch
$f'=af$ eindeutig bestimmt ist.

Für eine Funktion $g(x)=f(x)\mathrm e^{-ax}$ folgt mit $f'=af$, dass $g'=0$
ist; damit ist $g$ gleich einer Konstante und man kann $f$ daraus
zurückrechnen. Die Umkehrung — dass $f(x)=\lambda\mathrm e^{ax}$ tatsächlich
$f'=af$ erfüllt — klärt, obund jede solche Funktion wirklich so aussieht.

\textit{Skript:} Kap. 4, Mittelwertsatz bzw. Konstanz folgt aus $g'=0$.
\end{werkzeugkasten}

\begin{achtung}[Darauf achten]
\begin{itemize}
  \item Produktregel vollständig anwenden: Ableitung von $f(x)\mathrm e^{-ax}$
        ist $f'(x)\mathrm e^{-ax}-a f(x)\mathrm e^{-ax}$, nicht ohne den
        inneren Term.
  \item Vorzeichen: $-a$ bei $\mathrm e^{-ax}$; ein Vorzeichenfehler verändert
        das Ergebnis vollständig.
  \item Aus $g'=0$ folgt $g=\text{konst}$, nicht sofort ein bestimmter Wert;
        die Konstante bestimmt man über einen geeigneten Funktionswert (z. B.
        $x=0$).
\end{itemize}
\end{achtung}
